The module. Realise the six generators on $S = \C^{8}$, the Fock space of three ladder pairs of Proposition \claimref{correspondence.twistors-are-states-three-ladder-pairs}, by the Jordan-Wigner matrices $g_{1} = \sigma_{x}\otimes 1\otimes 1$, $g_{2} = \sigma_{y}\otimes 1\otimes 1$, $g_{3} = \sigma_{z}\otimes\sigma_{x}\otimes 1$, $g_{4} = \sigma_{z}\otimes\sigma_{y}\otimes 1$, $g_{5} = \sigma_{z}\otimes\sigma_{z}\otimes\sigma_{x}$, $g_{6} = \sigma_{z}\otimes\sigma_{z}\otimes\sigma_{y}$ (Hermitian, squares $1$, pairwise anticommuting), with $e_{0} = g_{1}$, $\nu = g_{2}$, $e_{1} = ig_{3}$, $e_{2} = ig_{4}$, $e_{3} = ig_{5}$, $e_{-} = ig_{6}$. The check verifies the Clifford relations of the stated signature, that the volume element is $i$ times the chirality $\sigma_{z}\otimes\sigma_{z}\otimes\sigma_{z}$, that every bivector commutes with the chirality, so $\CS$ and $\CSm$ (its two eigenspaces, the even and odd states of the pair system) are invariant and four-dimensional, and that $Z^{2} = -1$. The Hilbert form is the standard one on $\C^{8}$, and generators squaring to $\pm 1$ are Hermitian, respectively anti-Hermitian, as stated. Any other realisation of the module is unitarily equivalent to this one on each chirality half, so the statements below do not depend on the choice.

\emph{(i)} A sesquilinear form $\ip{\psi}{X\phi}$ is invariant under the rotor flow of $B$ exactly when $B^{*}X + XB = 0$. For $X = iZ$ this holds for all fifteen bivectors (computed); $iZ$ is Hermitian. Writing the fifteen conditions on $\CS$ as a linear system for the sixteen entries of a $4\times 4$ matrix $X$, exact elimination over the Gaussian rationals gives a kernel of complex dimension one, so every invariant sesquilinear form is a complex multiple of $h = \ip{\cdot}{iZ\,\cdot}$ and every invariant Hermitian form a real multiple. The same holds on $\CSm$. On $\CS$, in the Jordan-Wigner basis, $iZ = \mathrm{diag}(-1, -1, 1, 1)$ and $Z = \mathrm{diag}(i, i, -i, -i)$: signature $(2,2)$, $Z = -i$ on $T_{+}$ and $+i$ on $T_{-}$; $Z$ is the generator of the centre of the maximal compact subalgebra of Proposition \claimref{synthesis.positive-energy-selects-rule-2-sign}(i), whose identification with the attachment's energy that result left argued. The twistor form of the attachment is an invariant nondegenerate Hermitian form on $\CS$, hence a nonzero real multiple $c\,h$, and $h(\psi, iZ\phi) = \ip{\psi}{(iZ)^{2}\phi} = \ip{\psi}{\phi}$ since $(iZ)^{2} = -Z^{2} = 1$.

\emph{(ii)} $q_{B}(\psi) = i\ip{\psi}{iZB\psi}$ has matrix $i\,hB$ on $\CS$, Hermitian for every bivector $B$ by (i). For $B = Z$, $hZ = iZ\cdot Z = -i$, so $q_{Z} = \ip{\psi}{\psi}$. For the other generators the check computes $i\,hB$ exactly and tests all fifteen principal minors: positive semidefinite for $P_{a}$ with $a = e_{0}$, $e_{0} \pm e_{1}$, $e_{0} + \tfrac35 e_{2} + \tfrac45 e_{3}$, $e_{0} + e_{3}$ (null and timelike future vectors) and $2e_{0} + e_{1}$, $5e_{0} + 3e_{1} + 4e_{3}$ (timelike); the same for $a\nO$ and for $P_{a} + \tfrac{\Lambda}{12}a\nO$ at $\Lambda = 3, \tfrac37, 12$, which for null $a = n$ is the de Sitter generator $\tfrac12 nI$; negative semidefinite for the past vectors $-a$; indefinite for the spacelike $a = e_{1}$, $e_{2}$, $e_{0} + 2e_{1}$, and for the boost $e_{0}e_{1}$, the dilation $\nu e_{-}$ and the boost $e_{1}e_{-}$. For general future $a$ the statement follows from these by the structure of the cone: $a \mapsto q_{P_{a}}$ is linear, and the future cone is generated by future null vectors, all positive multiples of $e_{0} + \hat n$ with $\hat n$ a unit spatial vector, and so conjugate to $e_{0} + e_{1}$ up to positive scale under the spatial rotations, which commute with $Z$, while $q$ transforms by congruence under those rotations (they preserve $h$ and commute with $Z$), so semidefiniteness and its sign are preserved; $a\nO$ is handled the same way.

\emph{(iii)} With $T_{-}$ spanned by the first two basis vectors and $T_{+}$ by the last two, set $\hat Z^{\alpha} = b_{\alpha}^{*}$ on $T_{-}$ and $\hat Z^{\alpha} = a_{\alpha}$ on $T_{+}$; then $[\hat Z^{\alpha}, \hat Z^{\alpha\,*}] = -1$, resp. $+1$, which is $(h^{-1})^{\alpha\alpha}$. The check realises the four oscillators on polynomials (creation as multiplication, annihilation as differentiation, exact). From $[\hat Z^{\alpha\,*}, \hat Z^{\gamma}] = -(h^{-1})^{\gamma\alpha}$ one gets $[Q(M), \hat Z^{\gamma}] = -(M\hat Z)^{\gamma}$, hence $[Q(M), Q(N)] = Q([M, N])$ for the bivector generators; the check confirms the homomorphism exactly on three pairs and two test states. Since $hZ = -i$, $Q(\tfrac12 Z) = -\tfrac{i}{2}\sum_{\alpha}\hat Z^{\alpha\,*}\hat Z^{\alpha} = -\tfrac{i}{2}(b_{1}b_{1}^{*} + b_{2}b_{2}^{*} + a_{1}^{*}a_{1} + a_{2}^{*}a_{2}) = -\tfrac{i}{2}(N_{a} + N_{b} + 2) = -iE$; the check confirms this on all $625$ monomials of degree at most $4$ in each oscillator. The trace of $\tfrac12 Z$ on $\CS$ vanishes, so this ordering agrees with the Weyl ordering and no constant is free. The centre of $\mathfrak{u}(2,2)$, $M = \tfrac{i}{2}\,\id$, has Weyl-ordered quadratic $\tfrac{i}{2}\cdot\tfrac12\sum_{\alpha}h_{\alpha\alpha}(\hat Z^{\alpha\,*}\hat Z^{\alpha} + \hat Z^{\alpha}\hat Z^{\alpha\,*}) = \tfrac{i}{2}(N_{a} - N_{b}) = is$ (checked). On helicity $s \ge 0$ the states have $N_{a} = N_{b} + 2s$, so $E = N_{b} + s + 1 = |s| + 1 + n$ with $n = N_{b}$, and the number of monomials is $(N_{a} + 1)(N_{b} + 1) = (n + 2|s| + 1)(n + 1)$; symmetrically for $s < 0$. The check tabulates every $(s, E)$ with $|s| \le 2$ and $n + 2|s| \le 4$. These are the energies and multiplicities of Proposition \claimref{field-theory.twistor-attachment}, and the charge $s$ is the distinguished $\mathfrak{u}(1)$ charge of Proposition \claimref{hawking.helicity-lift-joints-itemised}, up to the orientation of that charge, which the energies do not see since they depend on $|s|$. In the other quantisation the commutators are $[\hat Z^{\alpha}, \hat Z^{\beta\,*}] = -(h^{-1})^{\alpha\beta}$, so it quantises $-h$; there $Q(\tfrac12 Z) = +iE$ (checked on three monomials), a generator bounded above. For future $a$, write the positive semidefinite $i\,hP_{a} = \sum_{k}\lambda_{k}v_{k}v_{k}^{*}$ with $\lambda_{k} \ge 0$; then $iQ(P_{a}) = \sum_{k}\lambda_{k}c_{k}^{*}c_{k}$, $c_{k} = \sum_{\beta}\overline{(v_{k})_{\beta}}\hat Z^{\beta}$, a sum of positive operators. The check confirms nonnegative Fock expectation values for four future null translations on five test states, a strictly positive value on some, and a negative value for the past translations; on the lowest state of the module ($E = 1$, $s = 0$) the expectation of $iQ(P_{e_{0} + e_{1}})$ is $1$.

\emph{(iv)} The oscillator representation of $\mathfrak{su}(2,2)$ on each helicity module integrates to a unitary irreducible representation of $SU(2,2) \cong \mathrm{Spin}^{+}(2,4)$, the ladder representations (Mack and Todorov; Mack; cited from memory); its generator $iQ(\tfrac12 Z) = E$ is bounded below by $1$, so $\tfrac12 Z$, and hence $Z$, lies in the cone $C_{\pi}$ of nonnegative generators in the sense of Proposition \claimref{synthesis.positive-conformal-energy-gives-positive-null-translations}. Proposition \claimref{field-theory.twistor-attachment} states that the tower's CAR is the second quantisation over this positive-energy one-particle space (via the Penrose transform, Fact \claimref{field-theory.imported-standard-field-theory}). The second-quantisation functor sends a unitary $U$ to $\Gamma(U)$ on Fock space, fixes the vacuum, and sends a nonnegative generator $A$ to $d\Gamma(A) \ge 0$ (cited from memory, standard); so the Fock vacuum's GNS representation carries a unitary representation of $\mathrm{Spin}^{+}(2,4)$ with $Z \in C_{\pi}$ and invariant GNS vector. By the stored positivity result every future translation, in particular every null translation along the generators of a local Rindler horizon of the flat background, has a nonnegative generator, and the vacuum is invariant under it.

\emph{(v)} The Fock vacuum is quasi-free (its correlations are those of the one-particle projection) and gauge-invariant (the number charge annihilates it). Its restriction to the field on the light rays of a null plane of Minkowski space is then a gauge-invariant quasi-free state invariant under the null translations of the plane, with nonnegative generator: the clause of Premise (N$_{+}$) in case (c), read with the horizon modes the series' field on those rays: its canonical component on the null plane $N_{k}$ is $\Lambda(m)\psi$, $m$ the dual null vector, by clause (i) of the result filed alongside (synthesis.boost-evolution-on-generators-is-one-unitary-group), on which the restriction is the continuum mode space of each generator (the light-ray algebra of a free massless field, cited from memory, Wall; Casini, Teste and Torroba). No temperature, period or Euclidean section enters. For the de Sitter field, (ii) gives the same sign for $\tfrac12 nI$, which together with Proposition \claimref{synthesis.de-sitter-horizon-translation-is-conformal-null-rotation}(iii) to (v) gives the clause in case (b) for the conformally invariant fields of the attachment. By Theorem \claimref{synthesis.null-vacuum-on-a-horizon-generator}(ii) a state satisfying the clause on a generator is unique (the filled sea), so on these cases the clause is equivalent to the statement that the state of the light rays is the attachment's vacuum restricted to them (in case (b), the conformal vacuum). This shows that the vacuum satisfies the clause and restates (N$_{+}$) there as the choice of that vacuum; it does not derive the choice. For case (a) the stored de Sitter result shows that the Schwarzschild field admits no gauge in which its horizon translations are background conformal transformations, so nothing here reaches it. $\square$
